%==============================================================================
%  CAPÍTULO — ARBITRAJE CONCURSAL
%==============================================================================
\chapter{El arbitraje concursal}
\label{cap:arbitraje}

\fichaCapitulo{La jurisdicción convencional del concurso.}{Capítulo VII de la ley:
sometimiento a arbitraje, nómina de árbitros concursales, quórums de decisión y reglas de
tramitación.}

%==============================================================================
\bloqueTeorico
%==============================================================================

\subsection{Un procedimiento colectivo ante jurisdicción convencional}

El arbitraje concursal plantea una tensión conceptual que conviene enunciar con
precisión. El concurso es un procedimiento universal y de orden público, cuyos efectos
alcanzan a acreedores que no han comparecido y a terceros que no han consentido. El
arbitraje, en cambio, descansa sobre la autonomía de la voluntad y produce efectos entre
quienes se sometieron a él.

La \LIR{} \citeyearpar{ley-20720} resuelve la tensión mediante un diseño particular: no
somete el concurso a arbitraje por acuerdo bilateral entre deudor y un acreedor, sino por
decisión de la junta de acreedores adoptada con los quórums que la ley
establece.\verificar{Precisar el quórum exigido y la oportunidad para adoptar la decisión
en el texto vigente del Capítulo VII.} La voluntad relevante no es la individual, sino la
colectiva del órgano que representa a la masa.

\subsection{Fundamento de la institución}

Las razones que explican la existencia de esta vía son de orden práctico:

\begin{enumerate}
  \item \textbf{Especialización.} Los concursos de gran complejidad ---grupos
  empresariales, activos de valoración técnica difícil, pasivos con estructuras
  financieras sofisticadas--- exigen un conocimiento que el juez civil de competencia
  común no siempre puede dedicar.

  \item \textbf{Celeridad.} El tribunal arbitral dedica al asunto una atención exclusiva
  de la que carece un juzgado con carga general.

  \item \textbf{Previsibilidad para el acreedor financiero.} La posibilidad de sustraer
  el concurso a la agenda judicial ordinaria reduce el descuento por incertidumbre con
  que los acreedores valoran su recuperación esperada.
\end{enumerate}

\begin{alerta}[El costo como límite real]
El arbitraje concursal es una vía cara: los honorarios del tribunal arbitral se solventan
con cargo a la masa. Su empleo sólo resulta racional cuando el activo concursal justifica
el gasto, lo que en la práctica lo reserva a concursos de envergadura.
\end{alerta}

%==============================================================================
\bloqueNormativo
%==============================================================================

\subsection{Ámbito de aplicación}

El arbitraje puede recaer tanto sobre el procedimiento de reorganización como sobre el de
liquidación. La decisión y su oportunidad difieren en cada caso, y son el primer punto
que debe verificarse.\verificar{Determinar, en el texto vigente, la oportunidad y el
órgano competente para acordar el arbitraje en cada procedimiento.}

\subsection{La nómina de árbitros concursales}

La ley encomienda a la \superir{} la formación y mantención de una nómina de árbitros
concursales, sujeta a requisitos de idoneidad y a un régimen de inhabilidades análogo al
de los demás auxiliares.\verificar{Confirmar si la nómina de árbitros subsiste tras la
Ley 21.563 y qué norma administrativa la regula.}

\subsection{Reglas de tramitación}

El tribunal arbitral aplica las normas de la ley concursal, con las adaptaciones propias
de la sede. Subsisten íntegramente el régimen de publicidad por el \boletin{}, los
deberes de información hacia la \superir{} y la fiscalización de los auxiliares.

\subsection{Costos y honorarios}

Los honorarios del árbitro y los gastos del procedimiento se solventan con cargo a la
masa y quedan sujetos al control de la cuenta final.\verificar{Precisar el régimen de
determinación y control de los honorarios arbitrales.}

%==============================================================================
\bloquePractico
%==============================================================================

\subsection{Cuándo conviene proponer el arbitraje}

\begin{foco}[Criterios de decisión]
El arbitraje concursal merece considerarse cuando concurren, al menos, tres de estas
condiciones: activo concursal de magnitud suficiente para absorber el costo; complejidad
técnica en la valoración o realización de los bienes; pluralidad de acreedores
financieros sofisticados y coordinados; conflictos previsibles que exigirán resolución
frecuente; y agenda judicial local especialmente congestionada.
\end{foco}

\subsection{Preparación de la propuesta a la junta}

\begin{checklist}[Antes de someter la decisión a la junta]
  \item Estimación del costo arbitral y comparación con el escenario judicial.
  \item Identificación de árbitros de la nómina sin inhabilidades respecto del deudor ni
  de los acreedores principales.
  \item Sondeo previo de los acreedores que concentran el pasivo con derecho a voto.
  \item Evaluación del efecto sobre el calendario del procedimiento.
  \item Análisis del régimen recursivo aplicable y de sus diferencias con la sede
  judicial.
\end{checklist}

\begin{practica}[El punto que suele omitirse]
La decisión de arbitraje altera el régimen de impugnación de las resoluciones. Antes de
proponerla a la junta debe estar claro qué recursos subsisten, ante quién y con qué
efecto: es información que los acreedores necesitan para votar informadamente, y su
omisión puede fundar una impugnación posterior del acuerdo.
\end{practica}

%==============================================================================
\bloqueJurisprudencial
%==============================================================================

\subsection{Aplicación práctica de la institución}

El arbitraje concursal es, en la experiencia chilena, una institución de uso
infrecuente. El análisis debe abordar las razones de esa infrautilización ---costo,
desconocimiento, inercia procesal--- y contrastarlas con la experiencia comparada.
\pendiente{Sistematizar los casos chilenos en que se ha empleado el arbitraje concursal y
las razones de su escasa aplicación.}

\subsection{Límites de la competencia arbitral}

\pendiente{Determinar qué materias quedan sustraídas al conocimiento del árbitro:
acciones revocatorias, incidente de mala fe, calificación penal de la insolvencia y
tercerías de terceros no comparecientes.}
